\documentclass[a4paper,10pt]{article}
\usepackage[utf8]{inputenc}
\usepackage[margin=1.5cm, bottom=2cm]{geometry}
\usepackage{fancyhdr}
\usepackage[dvipsnames]{xcolor}
\usepackage{enumitem}
\usepackage{xcolor}
\usepackage[table]{xcolor}
\usepackage{makeidx}
\usepackage[most]{tcolorbox}
\newtcolorbox{osservazioni}{
  colback=gray!8,    % sfondo grigio chiaro
  colframe=white,      % nessun bordo
  arc=2mm,            % angoli smussati
  boxrule=0pt,        % spessore bordo 0
  fontupper=\footnotesize,
  left=2mm, right=2mm, top=2mm, bottom=2mm,
  before skip=7pt,   % niente spazio prima
  after skip=15pt     % niente spazio dopo
}
\newtcolorbox{osservazioni2}{
  colback=gray!8,    % sfondo grigio chiaro
  colframe=white,      % nessun bordo
  arc=2mm,            % angoli smussati
  boxrule=0pt,        % spessore bordo 0
  fontupper=\small,
  left=5mm, right=5mm, top=3mm, bottom=3mm,
  before skip=0pt,   % niente spazio prima
  after skip=30pt     % niente spazio dopo
}
\pagestyle{fancy}
\usepackage{tikz}
\usetikzlibrary{automata, positioning}
\pagestyle{fancy}

\usepackage{makeidx}

%%%%%%%%%%%%%%%%%%%
\newcommand{\EsercitazioneNum}{11}
\newcommand{\Data}{28/11/2025}
\newcommand{\DocType}{Esercitazione} % o "Eserciziario"
\newcommand{\FullTitle}{\DocType\ \EsercitazioneNum}

\newcommand{\Header}{
\noindent
  \small \textbf{Computabilità, Complessità e Logica - a.a. 2025/26} \hfill \textbf{\FullTitle} \\[0.2cm]
  \scriptsize Matteo Liotta, \texttt{matteo.liotta@studenti.units.it} \hfill \Data \\[0.2cm]
  \noindent\makebox[\linewidth]{\rule{\linewidth}{0.4pt}} \\[0.3cm]}

\newcommand{\NewPageHeader}{
  \vfill
  \noindent\makebox[\linewidth]{\rule{\linewidth}{0.4pt}}
  \newpage
  {\Header}
}
%%%%%%%%%%%%%%%%%%%

% Numero di pagina in basso al centro
\fancyhf{}
\fancyfoot[C] \thepage
\renewcommand{\footrulewidth}{0pt} % niente linea in basso
\renewcommand{\headrulewidth}{0pt} % niente linea in alto

\begin{document}

\footnotesize  % font generale più piccolo
\Header\\[0.7cm]


%%%%%%%%%%%%%%%% FOGLIO 2

    
    %% INDEX
    \addcontentsline{toc}{section}{Esercitazione 11: \textit{CNF e Algoritmi di Soddisfacibilità}}
    
    \noindent\Large \textbf{Esercitazione 11}: \textit{CNF e Algoritmi di Soddisfacibilità (DP, DPLL)} 
    \begin{osservazioni}
        \textbf{\textit{Nota}}. Sono indicati in \textcolor{ForestGreen}{\textit{verde}} gli esercizi a sfondo teorico, proposti per rafforzare le idee alla base dei metodi utilizzati. Sono invece indicati in \textcolor{blue}{\textit{blu}} gli esercizi che non saranno trattati durante l'esercitazione, bensì lasciati come strumento aggiuntivo di esercitazione individuale. Resta comunque assoluta disponibilità per la loro risoluzione negli incontri successivi, qualora necessario. 
    \end{osservazioni}
    
    {\normalsize
    \begin{itemize}[label=, leftmargin=0pt]
    
    %\item \textbf{\underline{Formule booleane cont'd.}}

    %% INDEX
    %\addcontentsline{toc}{subsection}{Esercizio 4.}
    %%
    %\item \textbf{\textcolor{black}{Esercizio 4.} (**)}\\
    %Si considerino le seguenti formule booleane:    
    %    $$\Big(((A \rightarrow (B \land \lnot C)) \lor (\overline A \lor B)) \rightarrow (\lnot A \lor C)\Big) \lor (B \land \lnot A) $$
    %    \noindent e
    %    $$\phi_2 = \big((\lnot A \lor B) \land (\lnot C \lor B)\big) \rightarrow ((C \lor B) \land A)$$
    
    
    %Si indichi, giustificando il motivo, se la prima è equivalente alla seconda.\\[0.2cm]

    
    \item \textbf{\underline{Formule booleane in forma CNF}}
    
    %% INDEX
    \addcontentsline{toc}{subsection}{Esercizio 7.} 
    %%
    \item \textbf{\textcolor{ForestGreen}{Esercizio 7.} (**)}\\ 	
    Si consideri la seguente formula booleana:
    $$\phi:= l_1\lor l_2 \lor l_3 \lor l_4$$
    Considerata allora la scrittura:
    $$\phi':=\phi_1\land \phi_2 = (X\lor l_1\lor l_2)\land(\overline{X} \lor l_3\lor l_4)$$
    Data una assegnazione $\alpha$ tale che $\alpha \models \phi$, si indichino tutte le assegnazioni che modellano $\phi'$.


    %% INDEX
    \addcontentsline{toc}{subsection}{Esercizio 8.} 
    %%
    \item \textbf{\textcolor{blue}{Esercizio 8.} (**)}\\ 	
    Si consideri la seguente formula booleana:

    $$\phi:= \bigg(P \rightarrow (Q\lor R)\bigg)\rightarrow \bigg( (P\rightarrow Q) \lor (P\rightarrow Q) \bigg)$$
    
    Si scriva $\phi$ in forma congiuntiva di clausole CNF.

    \begin{osservazioni}
    \underline{\textit{Soluzione.}}\\
    L'esercizio richiede la riformulazione in forma normale congiuntiva della formula booleana indicata.

    \textbf{1.} Riscriviamo le implicazioni logiche
    
    $$\phi:= \bigg(P \rightarrow (Q\lor R)\bigg)\rightarrow \bigg( (P\rightarrow Q) \lor (P\rightarrow R) \bigg)=$$
    $$\phi:= \overline {\bigg(P \rightarrow (Q\lor R)\bigg)} \lor \bigg( (P\rightarrow Q) \lor (P\rightarrow R) \bigg)=$$
    $$\phi:= \overline {\bigg(\overline P \lor (Q\lor R)\bigg)} \lor \bigg( (\overline P\lor Q) \lor (\overline P\lor R) \bigg)=$$

    \textbf{2.} Portiamo le negazioni alle variabili (NNF)
    $$\phi:=  {\bigg( P \land \overline {(Q\lor R)}\bigg)} \lor \bigg( (\overline P\lor Q) \lor (\overline P\lor R) \bigg)={\bigg( P \land (\overline Q\land \overline R)}\bigg) \lor \bigg( (\overline P\lor Q) \lor (\overline P\lor R) \bigg)$$

    Allora 
    $$=\bigg( P \land \overline Q\land \overline R\bigg) \lor \bigg( \overline P\lor Q \lor R \bigg)=$$

    Per semplicità, indichiamo 

    $$=\bigg( P \land \overline Q\land \overline R\bigg) \lor \tilde \phi =$$
    $$=(\tilde \phi \lor P) \land (\tilde \phi \lor \overline Q) \land (\tilde \phi \lor \overline R)$$

    E reintroducendo $\tilde \phi$:

    $$=\bigg((P \land \overline Q\land \overline R) \lor P\bigg) \land \bigg((P \land \overline Q\land \overline R) \lor \overline Q\bigg) \land \bigg((P \land \overline Q\land \overline R) \lor \overline R\bigg)$$

    Dove, se

    $$\bigg((P \land \overline Q\land \overline R) \lor P\bigg)=(P\lor P) \land (P\lor \overline Q) \land (P \lor \overline R) = P \land (P\lor \overline Q) \land (P \lor \overline R)$$

    Allora 
    
    $$\phi = ... = %\bigg(P \land (P\lor \overline Q) \land (P \lor \overline R)\bigg) \land \bigg((\overline Q \lor P ) \land (\overline Q\lor \overline Q) \land (\overline Q \lor \overline R)\bigg) \land \bigg((\overline R \lor P ) \land (\overline R\lor \overline Q) \land (\overline R \lor \overline R)\bigg)=$$
    \bigg(P \land (P\lor \overline Q) \land (P \lor \overline R)\bigg) \land \bigg((\overline Q \lor P ) \land \overline Q\land (\overline Q \lor \overline R)\bigg) \land \bigg((\overline R \lor P ) \land (\overline R\lor \overline Q) \land \overline R\bigg)$$    
    
    \end{osservazioni}
    \NewPageHeader    

    \begin{osservazioni}
    Da cui:

    $$\phi = 
    P \land (P\lor \overline Q) \land (P \lor \overline R) \land (\overline Q \lor P ) \land \overline Q\land (\overline Q \lor \overline R) \land (\overline R \lor P ) \land (\overline R\lor \overline Q) \land \overline R=$$
    $$P \land \overline Q \land \overline R \land (P\lor \overline Q) \land (P\lor \overline R)\land (\overline Q \lor \overline R)$$    
    
    \end{osservazioni}
    %% INDEX
    \addcontentsline{toc}{subsection}{Esercizio 9.} 
    %%
    \item \textbf{\textcolor{black}{Esercizio 9.} (**)}\\ 	
    Si consideri la seguente formula booleana:

    $$\phi:=(\ \lnot(X_1\lor X_2)\ \lor \ (X_3 \lor X_4)\  )\ \land\ X_5$$
    
    Si scriva $\phi$ in forma congiuntiva di clausole CNF.

    \begin{osservazioni}
    \underline{\textit{Soluzione.}}\\
    $$\phi:=(\ \lnot(X_1\lor X_2)\ \lor \ (X_3 \lor X_4)\  )\ \land\ X_5=$$
    $$(\ (\overline X_1\land \overline X_2)\ \lor \ (X_3 \lor X_4)\  )\ \land\ X_5=$$
    $$(\ (\overline X_1\land \overline X_2)\ \lor \ \tilde \phi\  )\ \land\ X_5=$$
    $$(\ (\overline X_1 \lor \tilde \phi) \land (\overline X_2 \lor \tilde \phi\  ))\ \land\ X_5=$$
    $$(\overline X_1 \lor \tilde \phi) \land (\overline X_2 \lor \tilde \phi\ )\ \land\ X_5=$$
    $$(\overline X_1 \lor X_3 \lor X_4) \land (\overline X_2 \lor X_3 \lor X_4 \ )\ \land\ X_5$$
    \end{osservazioni}
    $\\[0.2cm]$

    \textbf{\underline{Algoritmi di Soddisfacibilità: \textit{Davis-Putnam} (DP)}}
    
    %% INDEX
    \addcontentsline{toc}{subsection}{Esercizio 10.} 
    %%
    \item \textbf{Esercizio 10 (*)}\\ 	
    Si consideri il seguente insieme di clausole:
    $$\{Q,\lnot R\}, \{R,\lnot S\},\{S,\lnot T\},\{T,\lnot U\},\{U,\lnot W\},\{W,\lnot Q\}$$
    
    \begin{enumerate}
        \item Si indichi se la formula booleana $\phi$ corrispondente è soddisfacibile utilizzando l'algoritmo di Davis-Putnam (DP);
        \item Se soddisfacibile, si proponga una valutazione delle  $\alpha$ tale per cui $\alpha\models\phi$.
    \end{enumerate}

    \begin{osservazioni}
    \underline{\textit{Soluzione}.}\\
    $$S_0 = \{Q,\lnot R\}, \{R,\lnot S\},\{S,\lnot T\},\{T,\lnot U\},\{U,\lnot W\},\{W,\lnot Q\}$$
    
    \textbf{1.} Pivot $R$
    \begin{itemize}
        \item \textbf{\textit{Clausole esonerate:}} $\{S,\lnot T\},\{T,\lnot U\},\{U,\lnot W\},\{W,\lnot Q\}$
        \item \textbf{\textit{Risolventi:}} $\{Q,\lnot S\}$
    \end{itemize}
    $$S_1 = \{S,\lnot T\},\{T,\lnot U\},\{U,\lnot W\},\{W,\lnot Q\}, \{Q,\lnot S\}$$
    
    \textbf{2.} Pivot $T$
    \begin{itemize}
        \item \textbf{\textit{Clausole esonerate:}} $\{U,\lnot W\},\{W,\lnot Q\}, \{Q,\lnot S\}$
        \item \textbf{\textit{Risolventi:}} $\{S,\lnot U\}$
    \end{itemize}
    $$S_2= \{U,\lnot W\},\{W,\lnot Q\}, \{Q,\lnot S\},\{S,\lnot U\}$$

    
    \textbf{3.} Pivot $W$
    \begin{itemize}
        \item \textbf{\textit{Clausole esonerate:}} $\{Q,\lnot S\},\{S,\lnot U\}$
        \item \textbf{\textit{Risolventi:}} $\{U,\lnot Q\}$
    \end{itemize}
    $$S_3 = \{Q,\lnot S\},\{S,\lnot U\},\{U,\lnot Q\}$$
    \end{osservazioni}
    \NewPageHeader

\begin{osservazioni}
    \textbf{4.} Pivot $S$
    \begin{itemize}
        \item \textbf{\textit{Clausole esonerate:}} $\{U,\lnot Q\}$
        \item \textbf{\textit{Risolventi:}} $\{Q,\lnot U\}$
    \end{itemize}
    $$S_4 = \{U,\lnot Q\},\{Q,\lnot U\}$$

    \textbf{5.} Pivot $U$
    \begin{itemize}
        \item \textbf{\textit{Clausole esonerate:}} $\emptyset$
        \item \textbf{\textit{Risolventi:}} $\{Q,\lnot Q\}$
    \end{itemize}
    $$S_5 = \{Q,\lnot Q\}$$

    E poiché vanno eliminate le tautologie, si ottiene $S_5=\emptyset$. Segue la soddisfacibilità del modello.\\

    \underline{\textit{Identificazione di un modello $\alpha \mid \alpha\models\phi$}}\\
    
    Possiamo considerare la costruzione del modello. La variabile $Q$ sarà assegnata a piacimento tra $F$ e $T$, dato il risultato tautologico.
    $$\alpha(Q)=False$$
    Se così, allora ritornando al passo precedente, ritroviamo che doveva essere soddisfatta $\{Q,\lnot U\}$. Da questo segue che necessariamente dovremo avere allora:
    $$\alpha(U)=False$$
    Allora, poiché avevamo $\{S,\lnot U\}$, segue per il medesimo motivo:
    $$\alpha(S)=False$$
    Analogamente per  $\{U,\lnot W\}$:
    $$\alpha(W)=False$$
    E per  $\{T,\lnot U\}$
    $$\alpha(T)=False$$
    Ed infine, $\{Q,\lnot R\}$ richiede:
    $$\alpha(R)=False$$

    $$\alpha := \{S\rightarrow False,\ T\rightarrow False,\ W\rightarrow False,\ U\rightarrow False,\ Q\rightarrow False\} $$ è tale che $\alpha \models \phi$.
    
\end{osservazioni}
    %% INDEX
    \addcontentsline{toc}{subsection}{Esercizio 11.} 
    %%
    \item \textbf{\textcolor{black}{Esercizio 11.} (**)}\\ 	
    Si consideri il seguente insieme di clausole:
    $$\{\lnot A,B\}, \{\lnot B,C\}, \{\lnot C,D\},\{\lnot D,A\},\{\lnot A, \lnot C\},\{B, C\}$$
    \begin{enumerate}
        \item Si indichi se la formula booleana $\phi$ corrispondente è soddisfacibile utilizzando l'algoritmo di Davis-Putnam (DP);
        \item Se soddisfacibile, si proponga una valutazione delle  $\alpha$ tale per cui $\alpha\models\phi$.\\[0.2cm]
    \end{enumerate}

\begin{osservazioni}
    \underline{\textit{Soluzione.}}\\
    $$S_0 = \{\lnot A,B\}, \{\lnot B,C\}, \{\lnot C,D\},\{\lnot D,A\},\{\lnot A, \lnot C\},\{B, C\}$$
    
    \textbf{1.} Pivot $A$
    \begin{itemize}
        \item \textbf{\textit{Clausole esonerate:}} $\{\lnot B,C\}, \{\lnot C,D\},\{B, C\}$
        \item \textbf{\textit{Risolventi:}} $\{B, \lnot D\}, \{\lnot C, \lnot D\}$
    \end{itemize}
    $$S_1 = \{\lnot B,C\}, \{\lnot C,D\},\{B, C\},\{B, \lnot D\}, \{\lnot C, \lnot D\}$$

    \textbf{2.} Pivot $B$
    \begin{itemize}
        \item \textbf{\textit{Clausole esonerate:}} $\{\lnot C,D\}, \{\lnot C, \lnot D\}$
        \item \textbf{\textit{Risolventi:}} $\{C\}, \{\lnot D,C\}$
    \end{itemize}
    $$S_2 = \{\lnot C,D\}, \{\lnot C, \lnot D\}, \{C\}, \{\lnot D,C\}$$

    \textbf{3.} Pivot $C$
    \begin{itemize}
        \item \textbf{\textit{Clausole esonerate:}} $\emptyset$
        \item \textbf{\textit{Risolventi:}} $\{D\},\{\lnot D\}, \{D,\lnot D\},\{\lnot D\}$
    \end{itemize}
\end{osservazioni}
\NewPageHeader
\begin{osservazioni}
    $$S_3 = \{D\},\{\lnot D\}, \{D,\lnot D\},\{\lnot D\}=\{D\},\{\lnot D\},\{\lnot D\}$$\\
    
    Da cui si osserva che si origina clausola vuota da $\{D\},\{\lnot D\}$. Ne segue l'insoddisfacibilità della formula booleana corrispondente.
\end{osservazioni}
    
    %% INDEX
    \addcontentsline{toc}{subsection}{Esercizio 12.} 
    %%
    \item \textbf{\textcolor{blue}{Esercizio 12.} (**)}\\ 	
    Si consideri il seguente insieme di clausole:
    $$\{A, B, C\}, \{\lnot A, \lnot B\}, \{\lnot B, \lnot C\}, \{A, \lnot C\}, \{\lnot A, D\}, \{B, \lnot D\}$$
    \begin{enumerate}
        \item Si indichi se la formula booleana $\phi$ corrispondente è soddisfacibile utilizzando l'algoritmo di Davis-Putnam (DP);
        \item Se soddisfacibile, si proponga una valutazione delle  $\alpha$ tale per cui $\alpha\models\phi$.\\[0.2cm]
    \end{enumerate}

    \begin{osservazioni}
    \underline{\textit{Soluzione.}}\\
    $$S_0 = \{A, B, C\}, \{\lnot A, \lnot B\}, \{\lnot B, \lnot C\}, \{A, \lnot C\}, \{\lnot A, D\}, \{B, \lnot D\}$$\\
    
    \textbf{1.} Pivot $C$ (per comodità rispetto ai risolventi)
    \begin{itemize}
        \item \textbf{\textit{Clausole esonerate:}} $\{\lnot A, \lnot B\}, \{\lnot A, D\}, \{B, \lnot D\}$
        \item \textbf{\textit{Risolventi:}} $\{A, B, \lnot B\}, \{A, B, A\}$
    \end{itemize}
    $$S_1 = \{\lnot A, \lnot B\}, \{\lnot A, D\}, \{B, \lnot D\},\{A, B, \lnot B\}, \{A, B, A\}=\{\lnot A, \lnot B\}, \{\lnot A, D\}, \{B, \lnot D\}, \{A, B\}$$

    \textbf{2.} Pivot $D$
    \begin{itemize}
        \item \textbf{\textit{Clausole esonerate:}} $\{\lnot A, \lnot B\}, \{A, B\}$
        \item \textbf{\textit{Risolventi:}} $ \{\lnot A, B\}$
    \end{itemize}
    $$S_2 = \{\lnot A, \lnot B\}, \{A, B\},\{\lnot A, B\}$$

    \textbf{3.} Pivot $B$
    \begin{itemize}
        \item \textbf{\textit{Clausole esonerate:}} $\emptyset$
        \item \textbf{\textit{Risolventi:}} $\{\lnot A, A\}, \{\lnot A\}$
    \end{itemize}
    $$S_2 = \{\lnot A, A\}, \{\lnot A\} =\{\lnot A\}$$
    
    Allora, essendo presente una clausola con un solo letterale, è possibile solo assegnare $$\alpha(A)=False$$ affinché sia $\phi$ soddisfatta. Se così, si conclude producendo l'insieme vuoto di clausole e ne segue la soddisfacibilità della formula $\phi$ corrispondente. \\
    Sapendo che $\alpha(A)=False$ e avendo clausola $\{A, \lnot C\}$ ne segue che $$\alpha(C)=False$$Avendo $\{A, B, C\}$ si necessita di $$\alpha(B)=True$$Vediamo inoltre che tutte le clausole sarebbero ora soddisfatte, indipendentemente dal valore di $D$. Allora, i modelli che soddisfano $\phi$ sono:
    \begin{itemize}
        \item $\alpha := \{A\rightarrow False,\ B\rightarrow True,\ C\rightarrow False,\ D\rightarrow False\} $
        \item $\alpha := \{A\rightarrow False,\ B\rightarrow True,\ C\rightarrow False,\ D\rightarrow True\} $
    \end{itemize}
    

    
    \end{osservazioni}

    %% INDEX
    \addcontentsline{toc}{subsection}{Esercizio 13.} 
    %%
    \item \textbf{\textcolor{blue}{Esercizio 13.} (**)}\\ 	
    Si consideri il seguente insieme di clausole:
    $$\{P, Q, S\}, \{P, Q, \lnot S\}, \{P, \lnot Q\}, \{\lnot P, R\}, \{\lnot P, \lnot R\}$$
    \begin{enumerate}
        \item Si indichi se la formula booleana $\phi$ corrispondente è soddisfacibile utilizzando l'algoritmo di Davis-Putnam (DP);
        \item Se soddisfacibile, si proponga una valutazione delle  $\alpha$ tale per cui $\alpha\models\phi$.\\[0.2cm]
    \end{enumerate}

    \NewPageHeader

    \begin{osservazioni}
    \underline{\textit{Soluzione.}}\\
    $$\{P, Q, S\}, \{P, Q, \lnot S\}, \{P, \lnot Q\}, \{\lnot P, R\}, \{\lnot P, \lnot R\}$$\\
    
    \textbf{1.} Pivot $R$
    \begin{itemize}
        \item \textbf{\textit{Clausole esonerate:}} $\{P, Q, S\}, \{P, Q, \lnot S\}, \{P, \lnot Q\}$
        \item \textbf{\textit{Risolventi:}} $\{\lnot P\}$
    \end{itemize}
    $$S_1 = \{P, Q, S\}, \{P, Q, \lnot S\}, \{P, \lnot Q\},\{\lnot P\}$$

    Questo impone necessariamente che 
    $$\alpha(P)=False$$
    E allora 

    $S_1\overset{\alpha(P)=F}{=}\{Q, S\}, \{Q, \lnot S\}, \{\lnot Q\}$

    Da cui necessariamente

    $$\alpha(D)=False$$

    E allora 

    $S_1\overset{\alpha}{=}\{S\}, \{\lnot S\}$

    Allora, scegliendo come pivot $S$, si ottiene la clausola vuota, da cui segue l'insoddisfacibilità della formula.
    
    \end{osservazioni}

    %% INDEX
    \addcontentsline{toc}{subsection}{Esercizio 14.} 
    %%
    \item \textbf{\textcolor{blue}{Esercizio 14.} (**)}\\ 	
    Si consideri il seguente insieme di clausole:
    $$\{A, B\}, \{\lnot A, C\}, \{\lnot B, D\}, \{\lnot C, \lnot D\}, \{C, \lnot B\}, \{A, \lnot D\}$$
    \begin{enumerate}
        \item Si indichi se la formula booleana $\phi$ corrispondente è soddisfacibile utilizzando l'algoritmo di Davis-Putnam (DP);
        \item Se soddisfacibile, si proponga una valutazione delle  $\alpha$ tale per cui $\alpha\models\phi$.\\[0.2cm]
    \end{enumerate}

    \begin{osservazioni}
    \underline{\textit{Soluzione.}}
    $$\{A, B\}, \{\lnot A, C\}, \{\lnot B, D\}, \{\lnot C, \lnot D\}, \{C, \lnot B\}, \{A, \lnot D\}$$\\
    
    \textbf{1.} Pivot $D$
    \begin{itemize}
        \item \textbf{\textit{Clausole esonerate:}} $\{A, B\}, \{\lnot A, C\}, \{C, \lnot B\}$
        \item \textbf{\textit{Risolventi:}} $\{\lnot B, \lnot C\}, \{\lnot B, A\}$
    \end{itemize}
    $$S_1 = \{A, B\}, \{\lnot A, C\}, \{C, \lnot B\},\{\lnot B, \lnot C\}, \{\lnot B, A\}$$

    \textbf{2.} Pivot $B$
    \begin{itemize}
        \item \textbf{\textit{Clausole esonerate:}} $\{\lnot A, C\}$
        \item \textbf{\textit{Risolventi:}} $\{A, C\},\{A, \lnot C\},\{A\}$
    \end{itemize}
    $$S_2 = \{\lnot A, C\},\{A, C\},\{A, \lnot C\},\{A\}$$
    
    Segue necessariamente che $\alpha(A)=True$

    $$S_2 \overset{\alpha}{=} \{C\}$$

    Ed allora anche $$\alpha(C)=False$$ e si ottiene l'insieme vuoto di clausole, da cui la soddisfacibilità.

    Data la clausola $\{C, \lnot B\}$ segue $$\alpha(B)=False$$ e dato che ora il valore di $D$ non altera il valore di verità della formula, due modelli sono per $\phi$:
    \begin{itemize}
        \item $\alpha := \{A\rightarrow True,\ B\rightarrow False,\ C\rightarrow False,\ D\rightarrow False\} $
        \item $\alpha := \{A\rightarrow True,\ B\rightarrow False,\ C\rightarrow False,\ D\rightarrow True\} $
    \end{itemize}    
    \end{osservazioni}
    
    \NewPageHeader
    \textbf{\underline{Algoritmi di Soddisfacibilità: \textit{Davis-Putnam-Logemann-Loveland} (DPLL)}}

    %% INDEX
    \addcontentsline{toc}{subsection}{Esercizio 15.} 
    %%
    \item \textbf{\textcolor{ForestGreen}{Esercizio 15.} (**)}\\ 	
    Si consideri la formula booleana in forma CNF:
    $$\phi = (P)\land(\overline P \lor R)\land (\overline R \lor Q)$$
    \begin{enumerate}
        \item Si indichi se la formula booleana $\phi$ corrispondente è soddisfacibile utilizzando l'algoritmo di Davis-Putnam-Logemann-Loveland (DPLL);
        \item Se soddisfacibile, si proponga una valutazione delle  $\alpha$ tale per cui $\alpha\models\phi$.
    \end{enumerate}

    \begin{osservazioni}
    \underline{\textit{Soluzione.}}\\
    \textcolor{ForestGreen}{\textbf{Idea}: \textit{L'esercizio mostra l'utilità legata ad avere formule in forma CNF.}}
    
    
    \end{osservazioni}
    

    %% INDEX
    \addcontentsline{toc}{subsection}{Esercizio 16.} 
    %%
    \item \textbf{Esercizio 16. (**)}\\ 	
    Si consideri il seguente insieme di clausole:
    $$\{A, B, \lnot C\}, \{\lnot A, C, D\}, \{\lnot B, C, \lnot D\}, \{\lnot C, D, E\}, \{\lnot D, \lnot E, B\}, \{\lnot A, \lnot B, E\}$$
    \begin{enumerate}
        \item Si indichi se la formula booleana $\phi$ corrispondente è soddisfacibile utilizzando l'algoritmo di Davis-Putnam-Logemann-Loveland (DPLL);
        \item Se soddisfacibile, si proponga una valutazione delle  $\alpha$ tale per cui $\alpha\models\phi$.
    \end{enumerate}

    \begin{osservazioni}
    \underline{\textit{Soluzione.}}\\
    $$\{A, B, \lnot C\}, \{\lnot A, C, D\}, \{\lnot B, C, \lnot D\}, \{\lnot C, D, E\}, \{\lnot D, \lnot E, B\}, \{\lnot A, \lnot B, E\}$$\\
    
    \begin{itemize}[label=*]
        \item Sia $\alpha(A)=True$. Allora 
        $$\{C, D\}, \{\lnot B, C, \lnot D\}, \{\lnot C, D, E\}, \{\lnot D, \lnot E, B\}, \{\lnot B, E\}$$
        \begin{itemize}[label=*]
            \item Sia $\alpha(B)=True$. Allora:
            $$\{C, D\}, \{C, \lnot D\}, \{\lnot C, D, E\}, \{E\}$$
            \begin{itemize}[label=*]
                \item Sia $\alpha(C)=True$. Allora:
                $$\{D,E\}, \{E\}$$            
                    $\quad$ * Sia $\alpha(D)=True$. Allora:
                    $$\{E\}$$
            \end{itemize}
        \end{itemize}
    \end{itemize}
    
    Da cui, con $\alpha(E)=True$ segue la soddisfacibilità, poiché esiste almeno un modello che soddisfi la formula booleana corrispondente. Infatti, con 

   $$\alpha := \{A\rightarrow True,\ B\rightarrow True,\ C\rightarrow True,\ D\rightarrow True, \ E\rightarrow True\} $$

   $$\alpha \models\phi$$
    
    \end{osservazioni}

    %% INDEX
    \addcontentsline{toc}{subsection}{Esercizio 17.} 
    %%
    \item \textbf{Esercizio 17. (**)}\\ 	
    Si consideri il seguente insieme di clausole:
    $$\{A, B\}, \{A, \neg B\}, \{\neg A, C\}, \{\neg A, \neg C\}$$
    \begin{enumerate}
        \item Si indichi se la formula booleana $\phi$ corrispondente è soddisfacibile utilizzando l'algoritmo di Davis-Putnam-Logemann-Loveland (DPLL);
        \item Se soddisfacibile, si proponga una valutazione delle  $\alpha$ tale per cui $\alpha\models\phi$.\\[0.2cm]
    \end{enumerate}

    \NewPageHeader

    \begin{osservazioni}
    \underline{\textit{Soluzione.}}\\
    $$\{A, B\}, \{A, \neg B\}, \{\neg A, C\}, \{\neg A, \neg C\}$$\\
    
    \begin{itemize}[label=*]
        \item Sia $\alpha(A)=True$. Allora: 
        $$\{C\}, \{\neg C\}$$
        Si ottiene la clausola vuota: è necessario effettuare \textit{backtrack}.
        \item Sia $\alpha(A)=False$. Allora:
        $$\{B\}, \{\neg B\}$$\\
        Si ottiene la clausola vuota: è necessario effettuare \textit{backtrack}.
    \end{itemize}
    Non è possibile effettuare nuovamente backtrack e non si è ottenuta l'insieme vuoto di clausole: segue l'insoddisfacibilità della formula booleana corrispodente. Idealmente, nessuna assegnazione di $A$ è permette a $\phi$ di essere verificata.
    \end{osservazioni}

    %% INDEX
    \addcontentsline{toc}{subsection}{Esercizio 18.} 
    %%
    \item \textbf{\textcolor{blue}{Esercizio 18.} (**)}\\ 	
    Si consideri il seguente insieme di clausole:
    $$\{P, Q, R\}, \{\lnot P, Q, S\}, \{\lnot Q, R, S\}, \{\lnot R, \lnot S, Q\}, \{\lnot P, \lnot Q, \lnot R\}, \{P, \lnot Q, \lnot S\}$$
    \begin{enumerate}
        \item Si indichi se la formula booleana $\phi$ corrispondente è soddisfacibile utilizzando l'algoritmo di Davis-Putnam-Logemann-Loveland (DPLL);
        \item Se soddisfacibile, si proponga una valutazione delle  $\alpha$ tale per cui $\alpha\models\phi$.
    \end{enumerate}

    \begin{osservazioni}
    \underline{\textit{Soluzione.}}\\
    
    $$\{P, Q, R\}, \{\lnot P, Q, S\}, \{\lnot Q, R, S\}, \{\lnot R, \lnot S, Q\}, \{\lnot P, \lnot Q, \lnot R\}, \{P, \lnot Q, \lnot S\}$$
    
    Potremmo sceglier qualsiasi letterale, ma concentriamoci su alcuni che semplifichino l'insieme di clausole, come $P$.
    \begin{itemize}[label=*]
        \item Sia $\alpha(P)=True$. Allora 
        $$\{Q, S\}, \{\lnot Q, R, S\}, \{\lnot R, \lnot S, Q\}, \{\lnot Q, \lnot R\}$$
        \begin{itemize}[label=*]
            \item Sia $\alpha(Q)=True$. Allora
            $$\{R, S\}, \{\lnot R\}$$
            Da qui necessariamente abbiamo che $\alpha(R)=False$. 
            $$\{S\}$$
            E allora necessariamente $$\alpha(S)=True$$
        \end{itemize}
    \end{itemize}

    Abbiamo trovato un modello che soffisfa $\phi$:

       $$\alpha := \{P\rightarrow True,\ Q\rightarrow True,\ R\rightarrow False,\ S\rightarrow True\} $$

    Ne segue dunque la soddisfacibilità della formula.
    
    
    \end{osservazioni}

    %% INDEX
    \addcontentsline{toc}{subsection}{Esercizio 19.} 
    %%
    \item \textbf{\textcolor{blue}{Esercizio 19.} (**)}\\ 	
    Si consideri il seguente insieme di clausole:
    $$\{X, Y, \lnot Z\}, \{\lnot X, Z, U\}, \{\lnot Y, Z, V\}, \{\lnot Z, \lnot U, W\}, \{\lnot V, \lnot W, Y\}, \{\lnot X, \lnot V, W\}, \{X, \lnot Y, V\}$$
    \begin{enumerate}
        \item Si indichi se la formula booleana $\phi$ corrispondente è soddisfacibile utilizzando l'algoritmo di Davis-Putnam-Logemann-Loveland (DPLL);
        \item Se soddisfacibile, si proponga una valutazione delle  $\alpha$ tale per cui $\alpha\models\phi$.
    \end{enumerate}

    \NewPageHeader

    \begin{osservazioni}
    \underline{\textit{Soluzione.}}\\
    
    $$\{X, Y, \lnot Z\}, \{\lnot X, Z, U\}, \{\lnot Y, Z, V\}, \{\lnot Z, \lnot U, W\}, \{\lnot V, \lnot W, Y\}, \{\lnot X, \lnot V, W\}, \{X, \lnot Y, V\}$$

    \begin{itemize}[label=*]
        \item Sia $\alpha(W)=True$. Allora:
        $$\{X, Y, \lnot Z\}, \{\lnot X, Z, U\}, \{\lnot Y, Z, V\}, \{\lnot V, Y\}, \{X, \lnot Y, V\}$$
        \begin{itemize}[label=*]
            \item Sia $\alpha(X)=True$
            $$\{Z, U\}, \{\lnot Y, Z, V\}, \{\lnot V, Y\}$$
            \begin{itemize}[label=*]
                \item Sia $\alpha(Y)=True$. Allora
                $$\{Z, U\}, \{Z, V\}$$
                Da cui sufficiente 
                $$\alpha(Z)=True$$
                $$\alpha(U)=True$$
                $$\alpha(V)=True$$
            \end{itemize}
        \end{itemize}
    \end{itemize}

    Abbiamo trovato un modello che soffisfa $\phi$:

    $$\alpha := \{W\rightarrow True,\ X\rightarrow True,\ Y\rightarrow True,\ Z\rightarrow True\, U\rightarrow True, V\rightarrow True\} $$

    Ne segue dunque la soddisfacibilità della formula.
    \end{osservazioni}

    
    %% INDEX
    \addcontentsline{toc}{subsection}{Esercizio 20.} 
    %%
    \item \textbf{\textcolor{blue}{Esercizio 20.} (**)}\\ 	
    Si consideri il seguente insieme di clausole:
    $$\{A, B, C, D, \overline E\}, \{C, \overline B\}, \{D, \overline A\}, \{B, \overline D\}, \{C, A, \overline E\}, \{A, B, C, D, E\},$$$$\{C, \overline D\}, \{\overline D, \overline A, B\}, \{E, \overline C, D\}, \{C, A, \overline E\}$$
    \begin{enumerate}
        \item Si indichi se la formula booleana $\phi$ corrispondente è soddisfacibile utilizzando \textbf{sia} l'algoritmo di Davis-Putnam-Logemann-Loveland (DPLL) che l'algoritmo di Davis-Putnam (DP);
        \item Se soddisfacibile, si proponga una valutazione delle  $\alpha$ tale per cui $\alpha\models\phi$.
    \end{enumerate}

    \begin{osservazioni}
    \underline{\textit{Soluzione.}}\\
    
    $$\{A, B, C, D, \overline E\}, \{C, \overline B\}, \{D, \overline A\}, \{B, \overline D\}, \{C, A, \overline E\}, \{A, B, C, D, E\},$$$$\{C, \overline D\}, \{\overline D, \overline A, B\}, \{E, \overline C, D\}, \{C, A, \overline E\}$$

    \begin{itemize}[label=*]
        \item Sia $\alpha(A)=True$. Allora:
        $$\{C, \overline B\}, \{D\}, \{B, \overline D\},\{C, \overline D\}, \{\overline D, B\}, \{E, \overline C, D\}$$
        Segue che necessariamente $$\alpha(D)=True$$
        $$\{C, \overline B\}, \{B\},\{C\}, \{B\}$$
        $$\{C, \overline B\}, \{B\},\{C\}$$
        Segue necessariamente
        $$\alpha(C)=True$$
        $$\alpha(B)=True$$
    \end{itemize}

    
    Abbiamo trovato un modello che soffisfa $\phi$:

    $$\alpha := \{A\rightarrow True,\ B\rightarrow True,\ C\rightarrow True,\ D\rightarrow True\} $$

    Ne segue dunque la soddisfacibilità della formula.
    
    \end{osservazioni}

\end{itemize}
\vfill
\noindent\makebox[\linewidth]{\rule{\linewidth}{0.4pt}}
\end{document}